\begin{proof}
The right side of \eqref{I:R:eq:fc-F1-root} increases continuously from zero as
$\tau$ increases from $q$, while $1-\tau$ decreases from $1-q$ to zero.
There is exactly one crossing in $(q,1)$. Equations
\eqref{I:R:eq:fc-general-threshold}--\eqref{I:R:eq:fc-scalar-Psi} identify it as the
projection optimum $R_{1,1}(q)$. The previously proved rank-one output
completion theorem supplies a real rank-one Hankel attaining this relaxation,
so its value is the exact real and complex F1 optimum.

For an algebraic comparison with the selected quartic, put
\[
 \mathcal L_q(\tau)=(1+q)(\tau-q)-(1-q)\tau^2,\qquad
 \mathcal A_q(\tau)=(1-q)^2\tau^2+(q^2+2q-1)(\tau-q).
\]
The radical equation implies
$\mathcal L_q(\tau)=2(\tau-q)\sqrt{q(1-\tau)}>0$, and direct expansion gives
\begin{align*}
 &\mathcal A_q(\tau)^2-4q^3(\tau-q)(1-2q+q\tau)\\
 &\qquad=(1-q)^2\bigl[\mathcal L_q(\tau)^2
                        -4q(1-\tau)(\tau-q)^2\bigr].
\end{align*}
Thus \eqref{I:R:eq:fc-F1-root} provides a monotone, unambiguous representation of
the already selected quartic branch.
\end{proof}
